\documentclass[12pt,a4paper]{article}

\usepackage[T1]{fontenc}
\usepackage[utf8]{inputenc}
\usepackage[romanian]{babel}
\usepackage{geometry}

\geometry{margin=2.5cm}

\title{Tutoriat 3\\Cautare Adversariala - Compromis intre timp si calitate}
\author{}
\date{}

\begin{document}

\maketitle

\section{Limitări date de resurse}

Deși algoritmul MiniMax produce decizii optime, utilizarea sa practică este limitată de resursele de calcul disponibile.

În jocurile complexe, numărul stărilor care trebuie explorate crește foarte rapid, ceea ce face imposibilă analiza completă a arborelui de joc.

\subsection{Explozia combinatorială}

Principala problemă a algoritmului MiniMax este explozia combinatorială.

La fiecare nivel al arborelui apar noi mutări posibile, iar numărul total de stări crește exponențial.

Dacă:

\begin{itemize} 
\item $b$ = factorul de ramificare; 
\item $d$ = adâncimea arborelui. 
\end{itemize}

atunci numărul aproximativ de noduri explorate este:

\[
b^d
\]

Această creștere exponențială reprezintă principala limitare a algoritmului.

\subsection{Exemplu numeric}

Presupunem:

\[
b = 10
\]

și

\[
d = 6
\]

Numărul de stări analizate devine:

\[
10^6 = 1\,000\,000
\]

Dacă adâncimea crește la:

\[
d = 10
\]

atunci:

\[
10^{10}
=
10\,000\,000\,000
\]

Numărul de poziții devine uriaș.

\subsection{Problema în șah}

În șah, factorul mediu de ramificare este aproximativ:

\[
b \approx 35
\]

Pentru o analiză de numai 8 mutări:

\[
35^8
\approx
2.25 \times 10^{12}
\]

Numărul de poziții este mult prea mare pentru a fi evaluat complet într-un timp rezonabil.

\subsection{Limitările de timp}

Calculul tuturor nodurilor poate necesita foarte mult timp.

În aplicațiile reale, deciziile trebuie luate rapid.

Exemple:

\begin{itemize}
\item un motor de șah trebuie să răspundă în câteva secunde;
\item un joc video trebuie să răspundă aproape instantaneu;
\item un sistem autonom trebuie să reacționeze în timp real.
\end{itemize}

Din acest motiv, explorarea completă a arborelui nu este practică.

\subsection{Limitările de memorie}

Pe lângă timpul de execuție, apare și problema memoriei.

Pe măsură ce arborele crește:

\begin{itemize}
\item trebuie generate mai multe stări;
\item trebuie memorate informații despre succesorii acestora;
\item costul de stocare crește.
\end{itemize}

Pentru jocuri complexe, memoria disponibilă poate deveni insuficientă.

\subsection{Compromisul între timp și calitate}

În practică există un compromis între:

\begin{itemize}
\item calitatea deciziei;
\item timpul de calcul.
\end{itemize}

O adâncime mai mare:

\begin{itemize}
\item oferă decizii mai bune;
\item necesită mai mult timp.
\end{itemize}

O adâncime mai mică:

\begin{itemize}
\item reduce timpul de calcul;
\item poate conduce la decizii mai puțin precise.
\end{itemize}

\subsection{Soluții utilizate în practică}

Pentru a depăși limitările de resurse se utilizează:

\begin{itemize}
\item limitarea adâncimii de căutare;
\item funcții de evaluare;
\item Alpha-Beta Pruning;
\item ordonarea mutărilor (Analizarea mai întâi a mutărilor promițătoare pentru a crește eficiența algoritmului Alpha-Beta);
\item tabele de transpoziție (Structuri de date care memorează pozițiile deja evaluate pentru a evita recalcularea acelorași stări.).
\end{itemize}

Aceste tehnici permit explorarea doar a celor mai promițătoare ramuri ale arborelui.

\subsection{Concluzie}

MiniMax oferă rezultate optime doar atunci când poate explora complet arborele de joc.

În aplicațiile reale acest lucru este rareori posibil din cauza limitărilor de timp și memorie.

Din acest motiv sunt necesare metode suplimentare care reduc numărul de stări analizate și permit luarea unor decizii bune într-un timp rezonabil.


















\section{Căutarea în adâncime limitată}

Din cauza dimensiunii foarte mari a arborilor de joc, explorarea completă folosind MiniMax este adesea imposibilă.

O soluție este limitarea adâncimii de căutare.

\subsection{Idee}

În loc să explorăm arborele până la stările terminale, ne oprim la o adâncime prestabilită.

\[
d = limita\ de\ adancime
\]

Nodurile aflate la această adâncime sunt evaluate folosind o funcție de evaluare.

\subsection{Exemplu}

Dacă limita este:

\[
d = 3
\]

algoritmul analizează doar primele trei niveluri ale arborelui.

Toate nodurile de la nivelul 3 sunt evaluate și utilizate de MiniMax pentru alegerea mutării.

\subsection{Avantaje}

\begin{itemize}
\item reduce timpul de execuție;
\item reduce numărul de stări analizate;
\item permite utilizarea MiniMax în jocuri complexe.
\end{itemize}

\subsection{Dezavantaje}

\begin{itemize}
\item nu mai garantează soluția optimă;
\item calitatea deciziei depinde de funcția de evaluare;
\item pot fi ignorate evenimente importante aflate mai adânc în arbore.
\end{itemize}

\subsection{Concluzie}

Căutarea în adâncime limitată reprezintă metoda principală prin care MiniMax poate fi utilizat în probleme reale, sacrificând o parte din acuratețe pentru a obține un timp de răspuns acceptabil.









\section{Funcții de evaluare}

În multe jocuri, explorarea completă a arborelui până la stările terminale este imposibilă.

Din acest motiv, MiniMax este oprit la o anumită adâncime, iar pozițiile rezultate sunt evaluate folosind o funcție de evaluare.

\subsection{Definiție}

O funcție de evaluare estimează cât de favorabilă este o stare pentru jucătorul MAX.

Aceasta asociază fiecărei stări o valoare numerică.

\[
Eval(s)
\]

unde $s$ reprezintă starea analizată.

\subsection{Interpretarea valorilor}

În general:

\begin{itemize}
\item valori pozitive → avantaj pentru MAX;
\item valori negative → avantaj pentru MIN;
\item valori apropiate de 0 → poziție echilibrată.
\end{itemize}

Exemplu:

\[
Eval(s)=50
\]

indică un avantaj pentru MAX.

\[
Eval(s)=-50
\]

indică un avantaj pentru MIN.

\subsection{Caracteristici}

O funcție de evaluare trebuie să fie:

\begin{itemize}
\item rapidă;
\item ușor de calculat;
\item corelată cu șansele reale de câștig;
\item capabilă să distingă pozițiile bune de cele slabe.
\end{itemize}

\subsection{Funcții de evaluare pentru șah}

În șah, o funcție de evaluare estimează avantajul unei poziții pe baza valorii pieselor.

Un exemplu simplu este:

\[
Eval(s)=9Q+5T+3N+3B+P
\]

unde:

\begin{itemize}
\item Q = numărul de regine;
\item T = numărul de ture;
\item N = numărul de cai;
\item B = numărul de nebuni;
\item P = numărul de pioni.
\end{itemize}

Scorul final se obține calculând diferența dintre piesele lui MAX și cele ale lui MIN.

\subsection{Funcții de evaluare pentru X și 0}

În jocul X și 0, o funcție de evaluare poate număra liniile care pot conduce la victorie.

Exemplu:

\[
Eval(s)
=
Linii_{MAX}
-
Linii_{MIN}
\]

unde:

\begin{itemize}
\item $Linii_{MAX}$ reprezintă liniile încă deschise pentru MAX;
\item $Linii_{MIN}$ reprezintă liniile încă deschise pentru MIN.
\end{itemize}

Cu cât valoarea este mai mare, cu atât poziția este mai favorabilă pentru MAX.

\subsection{Utilizare în MiniMax}

Atunci când se atinge limita de adâncime:

\[
d = d_{max}
\]

algoritmul nu mai generează succesori.

În schimb, aplică funcția de evaluare și utilizează valoarea obținută ca și cum ar fi valoarea unei stări terminale.

\subsection{Avantaje}

\begin{itemize}
\item reduc timpul de calcul;
\item permit analizarea jocurilor complexe;
\item oferă estimări rapide ale pozițiilor.
\end{itemize}

\subsection{Dezavantaje}

\begin{itemize}
\item reprezintă doar o aproximare;
\item pot evalua greșit anumite poziții;
\item calitatea rezultatului depinde de calitatea funcției de evaluare.
\end{itemize}

\subsection{Concluzie}

Funcțiile de evaluare permit utilizarea practică a algoritmului MiniMax prin înlocuirea evaluării exacte a pozițiilor cu estimări rapide ale calității unei stări.




\end{document}